\centerline{\bf \Huge Делимость и остатки}
\centerline{\bf \Huge Решения}

\problem{1}
На острове росло 3150 пальм. Пираты, приехавшие отдохнуть на остров, вырубают некоторое количество пальм. В тех краях появляются пираты только трёх шаек. Известно, что после прихода пиратов шайки «Длинная рука» количество пальм на острове уменьшается вдвое. Пираты шайки «Черный глаз» уменьшают количество пальм в пять раз, а после шайки «Острый зуб» остается только седьмая часть пальм. Через некоторое время на одном острове осталось всего 9 пальм. Докажите, что за это время в тех краях побывали пираты шайки «Острый зуб».

\textbf{Доказательство.} Пусть в некоторый момент на острове осталось $x$ пальм. Одна из шаек после посещения острова оставила после себя 1/n всех пальм, т.е. оставила $x / n$ пальм. Заметим, что $x / n$ - число целое, значит, число $x$ сократилось на множитель п. Начальное число пальм на острове - 3150 делится на 7, а конечное число пальм - 9 на 7 не делится. Значит, число 3150 сокрашали на 7, таким образом, шайка пиратов «Острый зуб» побывала на острове.

\problem{2}
Дед посадил 306 репок. Внучка, зайдя в огород, может выдернуть $1 / 2,1 / 5$, $1 / 8,1 / 11, \ldots$ всех имеющихся там репок. Жучка может выдернуть $1 / 3,1 / 6$, $1 / 9, \ldots$, а мышка $1 / 4,1 / 7,1 / 10, \ldots$ всех имеющихся к этому моменту в огороде репок. Вернувшись с покупками из города, дед обнаружил, что в огороде осталось только 6 репок. Заходила ли в огород Жучка?

\textbf{Ответ.} Заходила.

\textbf{Решение.} Пусть в некоторый момент было $x$ репок Если кто-то выкопал $1 / n$ этих х репок, то. осталось: $\frac{(n-1)}{n} x$ репок. Заметим, что остаться может только челое число репок. Таким образом, из числа х сократился множитель $n$ и появился множитель (n-1). Так как число 306 делится на 9, а число 6 делится только на 3, то из числа 306 должен был сократится множитель 3 (возможно не один). Но на такой множитель может пропасть только если $n$ кратно 3, т.е. после действий Жучки. Значит, Жучка на огород заходила.

\problem{3}
Курс акций компании "Рога и копыта" каждый день в полдень повышается или понижается на $1 \%$ (курс не округляется). Может ли курс акций дважды принять одно и то же значение?

\textbf{Ответ.} Нет.

\textbf{Решение.} Заметим, что при повышении курса акций он умножается на $\frac{101}{100}$, а при понижении -- на $\frac{99}{100}$. То есть если курс акций был равен $x$, то после $k$ повышений и $l$ понижений курс акций станет равным $\left(\frac{101}{100}\right)^k\left(\frac{99}{100}\right)^l x$. Если $x=\left(\frac{101}{100}\right)^k\left(\frac{99}{100}\right)^l x$, то $101^k 99^l=100^{k+l}$. Но в правой части этого равенства стоит четное число, а в левой -нечетное. Противоречие.

\problem{4}
а) Докажите, что число $\frac{1}{2}+\frac{1}{3}+\frac{1}{4}+\frac{1}{5}+\frac{1}{6}+\frac{1}{7}$ - не целое.

\textbf{Доказательство.} Приведем все дроби к общему знаменателю $2 \cdot 3 \cdot 4 \cdot 5 \cdot 6 \cdot 7=7$ ! (не наименьшему), тогда в числителе будет сумма шести слагаемых: $3 \cdot 4 \cdot 5 \cdot 6 \cdot 7+2 \cdot 4 \cdot 5 \cdot 6 \cdot 7+2 \cdot 3 \cdot 5 \cdot 6 \cdot 7+2 \cdot 3 \cdot 4 \cdot 6 \cdot 7+2 \cdot 3 \cdot 4 \cdot 5 \cdot 7+2 \cdot 3 \cdot 4 \cdot 5 \cdot 6= =\frac{7!}{2}+\frac{7!}{3}+\frac{7!}{4}+\frac{7!}{5}+\frac{7!}{6}+\frac{7!}{7}$ Заметим, что слагаемое $\frac{7!}{5}$, в отличие от остальных слагаемых, не делится на 5, значит, вся сумма или числитель не делится на 5 . Но знаменатель $7!$ на 5 делится, значит он не сможет сократится и число $\frac{1}{2}+\frac{1}{3}+\frac{1}{4}+\frac{1}{5}+\frac{1}{6}+\frac{1}{7}$ не может быть целым.

б) Докажите, что число $\frac{1}{2}+\frac{1}{3}+\frac{1}{4}+\ldots+\frac{1}{99}+\frac{1}{100}$ - не целое.

\textbf{Доказательство.} Аналогично 4 а) приведем все дроби к общему знаменателю 100! (не наименьшему), тогда в числителе будет сумма 99 слагаемых: $\frac{100!}{2}+\frac{100!}{3}+\ldots+\frac{100!}{99}+\frac{100!}{100}$ Заметим, что слагаемое $\frac{100!}{97}$, в отличие от остальных слагаемых, не делится на 97, значит, вся сумма или числитель не делится на 97 Но знаменатель 100 ! на 97 делится, значит, число $\frac{1}{2}+\frac{1}{3}+\frac{1}{4}+\frac{2}{2}+\frac{1}{99}+\frac{1}{100}$ не может быть целым.

\problem{5}
Камни лежат в трех кучках: в одной - 51 камень, в другой - 49 камней, а в третьей - 5 камней. Разрешается объединять любые кучки в одну, а также разделять кучку из четного количества камней на две равные. Можно ли получить 105 кучек по одному камню в каждой?

\textbf{Ответ.} Нельзя.

\textbf{Решение.} Если объединить все кучки в одну, то получим кучку со 105 камнями, с которой уже ничего сделать нельзя. Не объединять кучки тоже нельзя, так как ни одну из трех кучек нельзя разделить на две равные. Значит надо сначала объединить какие-то две кучки. Разберем три случая.
1) Если мы объединим кучки 49 и 51, то получится две кучки 100 и 5. То есть число камней в каждой кучке будет делится на 5. Докажем, что свойство сохранится и далее.
В самом деле если мы объединяем две кучки в каждой из которых число камней делится на 5 то и число камней в новой кучке тоже будет делится на 5. (алгебраически - если в одной было $5 x$, а в другой $5 y$, то в сумме будет $5 x+5 y=5(x+y)$ - делится на 5). Если мы разделяем кучку число камней в которой делится на 5 делим на 2 равные, то и в каждой половине числа камней будет делится на 5. (если 5 х делится на 2 , то $x$ делится на 2 , а значит $u^{5 \frac{x}{2}}$ делится на 5.)
Значит, у нас будут получатся только кучки в которых количество камней делится на 5, то есть кучек из 1 камня не получится.
2) Объединим кучки 5 и 51, получится две кучки 56 и 49. Теперь получилось, что число камней в каждой кучке делится на 7. Значит, аналогично пункту 1) у нас все время будут получаться кучки в которых количество камней будет кратно 7, то есть кучек из одного камня не будет.
3) Объединим кучки 5 и 49, получится две кучки 54 и 51. Теперь получилось, что число камней в каждой кучке делится на 3. Значит, аналогично пункту 1) у нас все время будут получаться кучки в которых количество камней будет кратно 3, то есть кучек из одного камня не будет.
Мы разобрали все случаи и в каждом из них доказали, что получить 105 кучек по одному камню - нельзя.

\textit{Замечание}. Можно сказать, что в этой задаче есть инвариант - например в первом случае инвариантом является то, что количество камней в каждой кучке делится на 5.

\problem{6}
В Петиной коллекции 500 пауков. Каждую неделю происходит следующее: либо Петина коллекция пополняется на 300 пауков, либо в коллекции погибают какие-то 198 пауков. Какое минимальное количество пауков может остаться в Петиной коллекции?

\textbf{Ответ.} Два паука.

\textbf{Решение.} Заметим, что числа, на которые повышается или понижается количество пауков
(300 и 198) кратны 3. Значит убрать всех пауков (то есть сделать так, чтобы количество пауков изменилось на 500), Петя не сможет, так как 500 не делится на 3. Аналогично Петя не сможет сделать так, что у него останется ровно 1 паук - для этого надо уменьшить число пауков на 499, а 499 на 3 не делится. (Можно конечно объяснить еще проще у Пети не могло получится нечетного числа пауков)
Покажем, как у Пети могли остаться ровно 2 паука. Сначала два раза убьем по 198 пауков - останется 104. Далее легко увеличить число Петиных пауков на 6 - надо два раза добавить по 300, а потом три раза убрать по 198. Будем так увеличивать число Петиных пауков, пока их не станет 200 (это возможно, так как 200-104=96 - делится на 6). А после этого снова убьем 198 пауков. Останется ровно 2.

\textit{Замечание} 1. Решении этой задачи состоит из двух частей - доказательства, что не может остаться меньше 2 пауков и примера как может остаться ровно два паука. Так всегда в задачах, где просят найти минимум или максимум какой-то величины, в задачах на «оченку + пример»

\textit{Замечание} 2. Конечно решение придумывается в другом порядке. Сначала надо просто попробовать: сколько пауков можно оставить. Легко оставить 104 паука, потом можно получить 404, отсюда получается 206 и 8 пауков, далее можно получить 308, 110, 410, 212, 14, 314, 116, 416, 218, 20... Видно что число пауков отличается на что-то кратное 6. Отсюда можно понять, что увеличивая на 6 можно получить 200 и отсюда ровно 2. Также видно, что 0 получить скорее всего не выйдет, так как разность тогда будет не кратна 6.